% 31/10/2002 at 01:30
% Cinem�tica do Ponto Material: Movimento Circular.
% Written by TORDAR 

\documentclass[12pt]{article}  

\pagestyle{empty}

\usepackage{graphicx}
\usepackage{pst-all}

%Let's go: 

\begin{document} 
\null\vskip4cm
\begin{center} 
\begin{pspicture}(-1.1,0)(1.1,1.1)
\psset{unit=8cm}
\psarc[linewidth=3pt,arcsepA=3pt,arcsepB=3pt,linestyle=dashed,dash=5mm 1mm](0,0){0.75}{45}{135}
\psarc[linewidth=3pt,arcsepA=3pt,arcsepB=3pt]{->}(0,0){0.25}{0}{45}
\psline[linewidth=3pt,linecolor=blue]{->}(-1,0)(1,0)
\psline[linewidth=3pt,linecolor=blue]{->}(0,0)(0,1)
\psline[linewidth=5pt,linecolor=red]{->}(0,0)(-0.5303,0.5303)
\psline[linewidth=5pt,linecolor=red]{->}(0,0)( 0,0.75)
\psline[linewidth=5pt,linecolor=red]{->}(0,0)( 0.5303,0.5303)
\rput{0}(1.05,0){\scalebox{1.5}{$X$}}
\rput{0}(0,1.05){\scalebox{1.5}{$Y$}}
\rput{0}(-0.4,0.25){\scalebox{1.5}{$\mathbf{r}(t_3)$}}
\rput{0}(-0.1,0.4){\scalebox{1.5}{$\mathbf{r}(t_2)$}}
\rput{0}(0.4,0.25){\scalebox{1.5}{$\mathbf{r}(t_1)$}}
\rput{0}(0.2,0.8){\scalebox{1.5}{$S$}}
\rput{0}(0.3,0.15){\scalebox{1.5}{$\theta$}}
\end{pspicture}
\end{center}
\end{document} 
